\documentclass[]{article}
\usepackage{graphicx} % Required for inserting images
\usepackage{amsmath}
\usepackage{amssymb}
\usepackage{amsfonts}
\usepackage{stmaryrd}
\usepackage{titlesec}
\usepackage[margin=1in]{geometry}

\setlength{\parindent}{0pt}

\newcommand{\msp}{\text{ }}
\newcommand{\ligneh}{\msp | \msp}
\newcommand{\si}{\text{ si }}
\newcommand{\et}{\text{ et }}
\newcommand{\ou}{\text{ ou }}
\newcommand{\ch}{\text{ch}}
\newcommand{\sh}{\text{sh}}
\newcommand{\rg}{\text{rg}}
\newcommand{\Vect}{\text{Vect}}
\newcommand{\dl}[2]{\text{DL}_{#1}(#2)}
\newcommand{\mat}{\text{Mat}}
\newcommand{\ind}{\perp\!\!\!\perp} 
\newcommand{\conv}{\text{ converge }}
\newcommand{\cov}{\text{Cov}}
\newcommand{\tr}{\text{tr}}

\newcommand{\dd}{\text{d}}
\newcommand{\dx}{\text{d}x}
\newcommand{\dy}{\text{d}y}
\newcommand{\dz}{\text{d}z}
\newcommand{\dr}{\text{d}r}
\newcommand{\dt}{\text{d}t}
\newcommand{\dtheta}{\text{d}\theta}
\newcommand{\dphi}{\text{d}\phi}
\newcommand{\repr}{|\mathcal{R}}
\newcommand{\grad}{\overrightarrow{{\text{grad}}} \;}
\newcommand{\cste}{\text{cste}}

\setlength{\emergencystretch}{3em}
\hbadness=10000

\title{Demonstration de cours semaine 3}
\author{Huang Ziwen}
\date{2026-2027 PC*}

\begin{document}

\section{Un nombre fini d'espaces propres sont en somme directe}

Soit $u \in \mathcal{L}(E)$, $(\lambda_i)_{i \in \nint{1}{p}}$ des valeurs propres deux à deux distinctes. \\

Alors \underline{les $(E_{\lambda_i}(u))_{i \in \nint{1}{p}}$ sont en somme directe}. \\

Par récurrence sur $p$. $\forall p \in \mathbb{N}^*, H_p$ : ``$\forall (\lambda_i)_{i \in \nint{1}{p}}$ des valeurs propres deux à deux distinctes, les $(E_{\lambda_i}(u))_{i \in \nint{1}{p}}$ sont en somme directe.''

\begin{itemize}
    \item Initialisation : pour $p = 1$, $E_{\lambda_1}$ est en somme directe, donc $H_1$. 
    \item Hérédité : Supposons $H_p$. \\
    
    Soit $(\lambda_i)_{i \in \nint{1}{p+1}}$ des valeurs propres deux à deux distinctes. \\
    Soit $(x_i)_{i \in \nint{1}{p+1}} \in \prod_{i}^{p+1}E_{\lambda_i}(u)$ tel que $\sum_{i=1}^{p+1} x_i = 0$. \\
    
    Montrons que $\forall i \in \nint{1}{p+1}, x_i = 0$. \\

    \[ u(\sum_{i=1}^{p+1} x_i) = \sum_{i=1}^{p+1} u(x_i) = \sum_{i=1}^{p+1} \lambda_i x_i = 0 \]

    D'où \begin{align}
        \sum_{i=1}^{p+1} \lambda_i x_i &= 0 \\
        \sum_{i=1}^{p+1} x_i &= 0
    \end{align}

    En faisant (1) - $\lambda_{p+1}$(2), on a:
    \[ \sum_{i=1}^{p} (\lambda_i - \lambda_{p+1})x_i = 0 \]

    Par hypothèse de récurrence, les $(E_{\lambda_i})_{i \in \nint{1}{p}}$ sont en somme directe. \\
    
    Alors $\forall i \in \nint{1}{p}$, $(\lambda_i - \lambda_{p+1})x_i = 0$. Or les valeurs propres sont deux à deux distinctes, donc $\forall i \in \nint{1}{p}$, $x_i = 0$. On en déduit que $x_{p+1} = 0$. \\

    Alors $\forall i \in \nint{1}{p+1}, x_i = 0$. On a donc $H_{p+1}$. 
\end{itemize}

\section{Polynôme caractéristique d'un induit et $\dim E_\lambda \leq \omega(\lambda)$}

Soit $u \in \mathcal{L}(E), F$ un ssev stable par $u$ de dimension $p$. \\

Alors \underline{$\chi_{u_F}$ divise $\chi_u$}. \\

Soit $\mathcal{B} = (e_i)_{i \in \nint{1}{n}}$ une base adapté à $F$ ($\mathcal{B}' = (e_i)_{i \in \nint{1}{p}}$ est une base de $F$). \\

\[ \mat_{\mathcal{B}}(u) = \begin{pmatrix}
    A & (*) \\
    (0) & C
\end{pmatrix} \]
avec $A = \mat_{\mathcal{B}'}(u_F)$

Soit $x \in \mathbb{K}$
\begin{align}
    \chi_u(x) &= \det(x \text{Id}_E - u) \\
    &= \begin{vmatrix}
        xI_{p} - A & (*) \\
        (0) & xI_{n - p} - C
    \end{vmatrix} \\
    &= \det (x I_p - A) \det(xI_{n-p} - C) \\
    &= \chi_A(x) \chi_A(c) \\
\end{align} 

D'où
\[ \chi_u = \chi_A \chi_C = \chi_{u_F} \chi_C  \]

\underline{$\dim E_\lambda \leq \omega(\lambda)$}.  \\

Soit $\lambda \in \text{Sp}(u)$. En notant $\omega(\lambda)$ la multiplicité de $\lambda$ dans $\chi_u$, $\chi_u = (X - \lambda)^{\omega(\lambda)}Q$ avec $Q$ un polynôme tel que $Q(\lambda) \neq 0$. \\

$E_{\lambda}(u)$ est $u$-stable, considerons $u_{\lambda}$ l'induit.

\[ \chi_{u_{\lambda}} = (X - \lambda)^{\dim(E_{\lambda}(u))} \]

Par ce qui precède, $ (X - \lambda)^{\dim(E_{\lambda}(u))} $ divise $(X - \lambda)^{\omega(\lambda)}Q$. Par définition de l'ordre, $\dim E_{\lambda}(u) \leq \omega(\lambda)$

\section{Caractérisation des endomorphismes diagonalisables}

Soit $E$ un K-ev de dimension finie, $u \in \mathcal{L}(E)$. \\

\underline{$ u \text{ est diagonalisable} \iff \exists P \in \mathbb{K}[x] \text{ scindés à racines simples tq } P(u) = 0 $} \\

\begin{itemize}
    \item $\boxed{\implies}$
    
    Soit $\mathcal{B}$ une base de diagonalisation de $u$, ($\lambda_1, \dots, \lambda_p$) les valeurs propres distincts de $u$. $\mat_{\mathcal{B}}(u) = \text{diag}(\lambda_1, \dots, \lambda_1, \lambda_2, \dots, \lambda_p, \dots, \lambda_p)$. \\

    Pour $P \in \mathbb{K}[X]$, $\mat_{\mathcal{B}}(P(u)) = P(\mat_{\mathcal{B}}(u)) = \text{diag}(P(\lambda_1), \dots, P(\lambda_p))$\\

    Donc $P = \prod_{i=1}^{p} (X-\lambda_i)$ convient - $\mat_{\mathcal{B}}(P(u)) = P(\mat_{\mathcal{B}}(u)) = \text{diag}(0, \dots, 0)$ alors $P(u) = 0$

    \item $\boxed{\impliedby}$
    
    Soit $P = \prod_{i=1}^p(X-\alpha_i) \in \mathbb{K}[X]$ un polynôme annulateur de $u$, avec les $(\alpha_i)_{i \in \nint{1}{p}}$ deux à deux distincts. \\
    
    Alors $\prod_{i=1}^p(u-\alpha_i \text{Id}_E) = 0$. \\

    Soit $(L_1, \dots, L_p)$ la famille des polynôme interpolateurs de Lagrange associée à $(\alpha_i, \dots, \alpha_p)$. \\

    Soit $i \in \nint{1}{p}$
    \begin{align*}
        L_i &= \frac{\prod_{k=1, k \neq i}(X - \alpha_k)}{\prod_{k=1, k \neq i}(\alpha_i - \alpha_k)} \\
        (X-\alpha_i)L_i &= \frac{\prod_{k=1}(X - \alpha_k)}{\prod_{k=1, k \neq i}(\alpha_i - \alpha_k)} = \frac{P}{\prod_{k=1, k \neq i}(\alpha_i - \alpha_k)}
    \end{align*}

    En $u$,
    \[(u - \alpha_i \text{Id}_E) \circ L_i (u) = \frac{P(u)}{\prod_{k=1, k \neq i}(\alpha_i - \alpha_k)} = 0\]

    Donc $\text{Im } L_i(u) \subset \ker (u-\alpha_i \text{Id}_E) = E_{\alpha_i}(u)$, ceci $\forall i \in \nint{1}{p}$ \\

    Or, par définition, $1 = \sum_{i=1}^{p} L_i$ d'où $\text{Id}_E = \sum_{i=1}^{p} L_i(u)$. En particulier, $\forall x \in E, x = \sum_{i=1}^{p} L_i(u)(x)$ avec $\forall i \in \nint{1}{p}, L_i(u)(x) \in E_{\alpha_i}(u)$. \\

    Donc $E \subset \sum_{i=1}^p E_{\alpha_i}(u)$. Avec certains $E_{\alpha_i}$ qui sont éventuellement réduit à 0 ($\alpha_i \not \in \text{Sp}(u)$), $E \subset \bigoplus_{i=1, \alpha_i \in \text{Sp}(u)}^p E_{\alpha_i}(u)$. \\
    
    Finalement, $E = \bigoplus_{i=1, \alpha_i \in \text{Sp}(u)}^p E_{\alpha_i}(u)$. $u$ est diagonalisable. 
\end{itemize}

Soit $F$ un ssev de $E$ stable par $u$ diagonalisable. Alors \underline{l'induit $u_F$ est diagonalisable.} \\

$u$ est diagonalisable, alors fixons $P$ un polynôme annulateur de $u$ scindé à racines simples. \\

$P(u) = 0$ et a fortiori, $P(u_F) = 0$. Alors $u_F$ est diagonalisable.  

\section{Théorème de trigonalisation}

Soit $u \in \mathcal{L}(E)$. \\

\underline{$\chi_u$ est scindé $\iff$ il existe une base $\mathcal{B}$ de $E$ telle que $\mat_{\mathcal{B}}(u)$ est triangulaire supérieure}. 

\begin{itemize}
    \item $\boxed{\impliedby}$
    \[ M = \mat_{\mathcal{B}}(u) = \begin{pmatrix}
        \alpha_1 & & (*) \\
        & \ddots & \\
        (0) & & \alpha_n
    \end{pmatrix} \]

    Alors $\chi_u = \chi_M = \prod_{k=1}^{n}(X - \alpha_k)$. $\chi_u$ est scindé. 
    
    \item $\boxed{\implies}$
    Par récurrence forte. $\forall n \in \mathbb{N}^*, H_n$ : ``Pour $E$ de dimension $n$, $u \in \mathcal{L}(E)$, si $\chi_u$ est scindé, alors il existe une base $\mathcal{B}$ telle que $\mat_{\mathcal{B}}(u)$ est triangulaire supérieure''. 

    \begin{itemize}
        \item Initialisation : $H_1$ est vraie. 
        \item Hérédité : Soit $n \in \mathbb{N}^*$. Supposons que $\forall k \in \nint{1}{n}, H_k$ est vraie. \\
        
        Soit $E$ de dimension $n+1$.\\
        Soit $u \in \mathcal{L}(E)$ tel que $\chi_u = \prod_{i=1}^p (X - \alpha_p)$ est scindé. \\
        
        $\text{Sp}(u) \neq \emptyset$, alors fixons $\lambda \in \text{Sp}(u)$.

        \begin{itemize}
            \item Si $u = \lambda \text{Id}_E$, la matrice de $u$ dans toute base de $E$ est $\lambda I_{n+1}$ donc a fortiori triangulaire supérieure. 
            \item Sinon, posons $v = u - \lambda \text{Id}_E \neq 0$. \\
            
            En particulier, $\chi_v = \prod_{i=1}^p (X - \alpha_p + \lambda)$ est scindé. \\

            $\text{Im}(v)$ est un ssev stable par $v$ non réduit à $\{0\}$. Notons $w = v_{\text{Im}(v)}$. \\
            
            $\chi_w$ divise $\chi_v$ et $\chi_v$ est scindé, donc $\chi_w$ est scindé. \\

            Or, $v \neq 0$ donc $\text{Im}(v) \neq \{0\}$. De plus, $v$ n'est pas injective car $\lambda$ est une valeur propre, et $v$ est un endomorphisme en dimension finie, alors $v$ n'est pas surjective, $\text{Im} (v) \subsetneq E$. On a alors $1 \leq \dim(\text{Im}(v)) \leq n$. \\

            Par hypothèse de récurrence, fixon $\mathcal{B}' = (e_1, \dots, e_r)$ une base de $\text{Im}(u)$ telle que $\mat_{\mathcal{B}'}(u) $ est triangulaire supérieure. \\

            Complétons cette base en base $\mathcal{B} = (e_1, \dots, e_r, e_{r+1}, \dots , e_{n+1})$ de $E$. \\

            $ \mat_{\mathcal{B}}(v) = \begin{pmatrix}
                \mat_{\mathcal{B}'}(u) & (*) \\
                (0) & (0)
            \end{pmatrix} $ est triangulaire supérieure. 
        \end{itemize}

        Donc $H_{n+1}$, 
    \end{itemize}
\end{itemize}

\end{document}